\section{The Experimental Method}
\label{app:method}

The theory is checked on a bench, a finite, discrete, reproducible simulator that turns each question into a concrete test. This appendix sets down the standard every test is held to, because the honesty of the program lives in the method rather than in any one result.

Every claim is a single experiment. It builds the exact substrate, steps the knit over it, measures one thing, and returns a structured verdict, a status, the metrics it measured, the control it measured against, and the claim it supports or refutes. One experiment is one claim, and the build fails only on a code crash or a broken self-consistency check, never on an honest scientific negative, because a model that punished its own negatives would learn to stop reporting them.

Each verdict carries a depth grade, the rubric that keeps a passing test from being mistaken for a deep one.

\begin{table}[ht]
\centering
\begin{tabular}{@{}cp{0.62\textwidth}@{}}
\toprule
grade & what it establishes \\
\midrule
L0 & circular, the answer was put in by hand, kept only as an honest consistency note \\
L1 & a known mathematical fact correctly reproduced on the substrate \\
L2 & a known physical construction reproduced on the substrate \\
L3 & an emergent and novel result, with a control that could have failed \\
\bottomrule
\end{tabular}
\caption{The depth rubric. Only L3 is both new and guarded by a control, and an L3 claim presented without a control is automatically downgraded to partial.}
\label{tab:rubric}
\end{table}

Alongside the depth grade, a result carries a status tag that says how it was obtained, derived by argument or exact computation, measured in simulation against a control, predicted as a consequence not yet checked, or open. Depth says how strong the kind of result is, the tag says how it was established, and the two together are what the scoreboard reports. The discipline binding all of it is that a result counts as solid only when its control could have come out the other way, so a positive finding means something precisely because there was a way for it to fail.

Two further commitments keep the bench honest. The model is deterministic, never reaching for a random seed, and where a trend is wanted it is found by growing the size of a run rather than by averaging over draws, with real numbers appearing only as measured outputs and never inside the base. And every result is held to faithfulness, the requirement that it trace back to the eight elements with no ninth ingredient smuggled in, so a phenomenon that does not emerge from the eight is reported as a negative rather than forced by a quiet addition.

Measuring a hyperbolic bulk has its own hazards, because a finite patch is almost all boundary, and the bench meets them with dedicated tools, Bethe recursion and closed manifolds and band theory and kernel-polynomial methods, that read bulk-local observables rather than finite-patch artifacts. And the results are sorted into three layers by what they depend on, bulk findings that hold on any hyperbolic substrate, boundary findings that need the negative curvature, and flat-layer findings that depend on the cubic cusp being clean periodic space, with only the last sensitive to the cusp's periodicity, which a clean cusp passes and a disordered slice fails. This three-layer audit is how the paper knows which of its claims are robust and which still lean on the particular shape of the cusp.
